\newpage
\begin{center}
    \LARGE
    \textbf{\thesistitle}
    
    \vspace{0.5cm}
    
    \authornamefl
    
    \vspace{1cm}
    
    \textbf{Summary}
    
    \vspace{1cm}
\end{center}

Modern fitness and nutrition software applications heavily rely on manual calorie tracking 
and rigid, template-based meal databases. 
These traditional tools lack automated, mathematically precise personalization, requiring
 users to manually balance macro-nutrients through tedious trial and error. 

This thesis introduces FitTrack, an intelligent nutrition management platform designed to 
automate custom meal plan generation based on exact physiological requirements. 
The system implements an automated calculation engine driven by Linear Programming to 
dynamically scale individual ingredient weights down to the gram, ensuring all generated 
menus hit target caloric, protein, carbohydrate, and fat boundaries flawlessly. To 
complement the optimization engine, a Multi-Layer Perceptron (MLP) classifier is 
utilized to evaluate physiological metrics—such as age, sex, weight and height-to select 
the optimal baseline options from a predefined batch of candidate meals that best match 
the individual's user profile.

The application bridges the gap between passive food logging diaries and active, automated
nutritional coaching. Experimental evaluations demonstrate that the constraints 
optimization engine reliably satisfies all mathematical constraints within seconds, 
serving as a practical foundation for building personalized health and wellness 
applications.

\vspace{0.5cm}
\noindent \textbf{Keywords:} Automated Meal Planning, Linear Programming, Constraints Optimization, Multi-Layer Perceptron, Personalized Nutrition.